
\documentclass{article}
\usepackage{colortbl}
\usepackage{makecell}
\usepackage{multirow}
\usepackage{supertabular}

\begin{document}

\newcounter{utterance}

\centering \large Interaction Transcript for game `ifeval', experiment `ifeval', episode 59 with Qwen3.5{-}27B.
\vspace{24pt}

{ \footnotesize  \setcounter{utterance}{1}
\setlength{\tabcolsep}{0pt}
\begin{supertabular}{c@{$\;$}|p{.15\linewidth}@{}p{.15\linewidth}p{.15\linewidth}p{.15\linewidth}p{.15\linewidth}p{.15\linewidth}}
    \# & \multicolumn{2}{c}{Player} && \multicolumn{2}{c}{Game Master} \\
    \hline

    \theutterance \stepcounter{utterance}  
    & & & \multicolumn{4}{p{0.6\linewidth}}{
        \cellcolor[rgb]{0.9,0.9,0.9}{
            \makecell[{{p{\linewidth}}}]{
                \texttt{\tiny{[P1$\langle$GM]}}
                \texttt{Write a short proposal for a new research project that investigates how language evolves over time. I want to make it challenging, so:} \\
\texttt{1. Do not include any commas in your response.} \\
\texttt{2. Do not include the letter "c" anywhere in your response.} \\
\texttt{3. Your response should contain at least 250 words.} \\
            }
        }
    }
     \\ \\

    \theutterance \stepcounter{utterance}  
    & \multicolumn{4}{p{0.6\linewidth}}{
        \cellcolor[rgb]{0.561,0.737,0.561}{
            \makecell[{{p{\linewidth}}}]{
                \texttt{\tiny{[P1$\rangle$GM]}}
                \texttt{This plan aims to study how human speech shifts across many years. We will look at old texts and modern talk to find patterns. Our goal is to see why words change shape or meaning. Many factors drive this shift. People move to new lands. They meet other groups. They learn new things. All these things alter how we speak. We need to track these changes with great care. Our team will gather data from many sources. We will use old books and digital logs. We will also record live speech from many towns. This mix of data will help us build a strong model. We hope to find rules that explain why some words fade while others grow. Some terms vanish fast. Others stay for ages. We want to know what makes a word last. Is it simple sound? Is it deep meaning? Is it just luck? Our study will test these ideas. We will run many tests on huge sets of data. We will look for links between history and speech. Wars often change talk. Trade routes spread new terms. Technology brings fresh words too. We must see how all these forces work together. Our work will help us understand human thought better. Language shapes how we see the world. If words shift then our view shifts too. This project will show that link clearly. We need funds to buy tools and pay staff. We need time to sort through vast amounts of text. We need space to store all our findings. Please support this bold effort to map speech history. It is a hard task but very worth doing. We will learn much about our past and future. Our results will help teachers and writers too. They can see how talk grows and shrinks. This knowledge is key for all who study human life. We ask for your help to start this grand journey now. Let us explore the deep roots of speech together.} \\
            }
        }
    }
    & & \\ \\

    \theutterance \stepcounter{utterance}  
    & & & \multicolumn{2}{p{0.3\linewidth}}{
        \cellcolor[rgb]{0.9,0.9,0.9}{
            \makecell[{{p{\linewidth}}}]{
                \texttt{\tiny{[GM$|$GM]}}
                {'punctuation:no_comma': {}, 'keywords:letter_frequency': {'let_relation': 'less than', 'letter': 'c', 'let_frequency': 1}, 'length_constraints:number_words': {'relation': 'at least', 'num_words': 250}}
            }
        }
    }
    & & \\ \\

    \theutterance \stepcounter{utterance}  
    & & & \multicolumn{2}{p{0.3\linewidth}}{
        \cellcolor[rgb]{0.9,0.9,0.9}{
            \makecell[{{p{\linewidth}}}]{
                \texttt{\tiny{[GM$|$GM]}}
                \texttt{game\_result = LOSE} \\
            }
        }
    }
    & & \\ \\

\end{supertabular}
}

\end{document}
